% !TeX root = ../../Book4.tex
% 纲要标题：## 第9章　导出函子、Ext 与 Tor
% 纲要标签：b4:ch:08
% 纲要位置：计划纲要.md，第 3400 行。

\chapter{导出函子、Ext 与 Tor}
\label{b4:ch:08}
\BookChapterNumber{b4:ch:08}

\begin{chapterabstract}
本章将 Hom 和张量积不保持正合的缺口组织为 Ext 与 Tor，建立导出函子的消解计算、长正合列、消去性与维数平移，并给出典型模上的计算。
\end{chapterabstract}

\begin{learninggoals}
\begin{itemize}
\item 写出左、右导出函子的次数约定，并说明零次为何需要相应单侧正合性。
\item 通过消去性构造泛 \(\delta\)-函子的延拓，区分上同调与同调泛性的方向。
\item 在两个变量中计算 Ext 和 Tor，证明平衡性并正确使用左右模约定。
\item 从长正合列解释 Hom、张量积的正合性失败，使用一次消失判据。
\item 定义相对于消解的合冲模，推导正次数平移及零次的核、余核公式。
\end{itemize}
\end{learninggoals}

取整数 \(n\ge2\)，对短正合列 \(0\rightarrow\mathbb Z\xrightarrow{n}\mathbb Z\rightarrow\mathbb Z/n\mathbb Z\rightarrow0\) 施加 \(\operatorname{Hom}_{\mathbb Z}(-,\mathbb Z)\)。由于从有限群到 \(\mathbb Z\) 的同态只能为零，所得序列在自然识别下为
\[
0\longrightarrow0\longrightarrow\mathbb Z
\xrightarrow{n}\mathbb Z.
\]
Hom 的左正合性保证前三个位置正合，却不能使最后的乘 \(n\) 映射满射；其余核为 \(\mathbb Z/n\mathbb Z\)。投射消解把这一缺口识别为 \(\operatorname{Ext}_{\mathbb Z}^1(\mathbb Z/n\mathbb Z,\mathbb Z)\)，同时要求我们证明它不随消解的选择而变。张量积丢失的正合性也可用消解测量，这便得到 Tor。

\ProseParagraphBreak
本章需要\chapref{b4:ch:07}的比较定理和马蹄引理：比较映射的存在及同伦唯一性使消解定义的群与态射不依赖选择，马蹄引理的逐次分裂短正合列则在施加加性函子后仍正合，从而产生长正合列。\secref{b4:sec:0604}的有限对角双复形比较把按两个变量作消解的结果联系起来，用于证明 Ext 和 Tor 的平衡性。第二卷第6.1—6.3节给出张量积与 Hom 的基本性质，第7.3节证明投射模平坦。以下讨论 Ext 时两个变量都是同一含幺环的左模；讨论 Tor 时第一变量为右模，第二变量为左模。除明确写出的交换情形外，不增加底环交换性的假设。

\ProseParagraphBreak
导出函子也适用于模范畴以外的交换范畴。例如，对拓扑空间 \(X\) 上的交换群层取全局截面是左正合函子，其右导出给出层上同调。层上同调的具体计算见\secref{b4:sec:C03}。

\ProseParagraphBreak
本章可与 Weibel 的《An Introduction to Homological Algebra》\cite[Section 2.1, Definition 2.1.4; Sections 2.4--2.7]{Weibel1994HomologicalAlgebra}对照。左、右导出函子的泛性分别见\cite[Theorem 2.4.7; Section 2.5.1]{Weibel1994HomologicalAlgebra}，消去性的一般判据见\cite[Exercise 2.4.5]{Weibel1994HomologicalAlgebra}；Tor 与 Ext 的平衡性分别见\cite[Theorem 2.7.2, Theorem 2.7.6]{Weibel1994HomologicalAlgebra}。关于整数模的 Tor 与平坦性，另见\cite[Sections 3.1--3.2]{Weibel1994HomologicalAlgebra}。本章的双复形证明在每个总次数只需有限次消元。

% !TeX root = ../../Book4.tex
% 纲要标题：### 9.1 导出函子
% 纲要标签：b4:sec:0801
% 纲要位置：计划纲要.md，第 2650 行。

\section{导出函子}
\label{b4:sec:0801}

左正合或右正合函子丢失的信息，可以在施加函子后的消解同调中读出。比较定理保证所得对象不依赖消解的选择，并给出它们在态射上的作用；马蹄引理给出连接态射和长正合列；消去性进一步说明这套导出函子具有怎样的泛性质。


% 纲要内容（仅作写作计划，尚非教材正文）：
%
% * 左、右导出及次数约定
% * δ-函子与泛性质
% * 消去性与泛 δ-函子
%

\subsection{左、右导出及次数约定}
\label{b4:subsec:0801:01}

一个左正合函子已经正确保留了核，却可能把满态射变成非满态射。内射消解把这种失败放入一个复形，正次数上同调便记录原函子没有保留的信息。投射消解对于右正合函子起相反方向的作用。

\ProseParagraphBreak
例如，固定整数 \(n\ge2\)，对交换群取左正合函子 \(F=\operatorname{Hom}_{\mathbb Z}(\mathbb Z/n,-)\)。\cref{b4:exm:integer-injective-resolution}已经给出内射消解
\[
 0\longrightarrow\mathbb Z\longrightarrow\mathbb Q
 \longrightarrow\mathbb Q/\mathbb Z\longrightarrow0.
\]
施加 \(F\) 后，\(F(\mathbb Q)=0\)，而 \(F(\mathbb Q/\mathbb Z)\) 是 \(\mathbb Q/\mathbb Z\) 中被 \(n\) 零化的子群，即 \((\tfrac1n\mathbb Z)/\mathbb Z\cong\mathbb Z/n\)。所得复形为 \(0\to\mathbb Z/n\)，位于零次和一次。因此原函子在 \(\mathbb Z\) 上取值为零，消解却在一次留下 \(\mathbb Z/n\)。这正是 \(\mathbb Q\to\mathbb Q/\mathbb Z\) 经过 \(F\) 后失去满射性的部分；下面将这种计算定义为导出函子。

\begin{definition}\label{b4:def:derived-functors}
设 \(\mathcal A,\mathcal B\) 是交换范畴，\(F:\mathcal A\to\mathcal B\) 是加性函子。若 \(\mathcal A\) 有足够多内射对象，对每个对象 \(A\in\mathcal A\) 选内射消解 \(A\to I_A^\bullet\)，定义
\[
R^nF(A)\coloneqq H^n(F(I_A^\bullet)),\qquad n\ge0,
\]
称为 \(F\) 的第 \(n\) 个\term[youdaochuhanzi]{右导出函子}[right derived functor]。若 \(\mathcal A\) 有足够多投射对象，对每个对象 \(A\in\mathcal A\) 选投射消解 \(P^A_\bullet\to A\)，定义
\[
L_nF(A)\coloneqq H_n(F(P^A_\bullet)),\qquad n\ge0,
\]
称为第 \(n\) 个\term[zuodaochuhanzi]{左导出函子}[left derived functor]。对态射 \(f:A\to A'\)，选择延拓或提升 \(f\) 的消解比较映射，逐次应用 \(F\) 后取相应次数的同调，所得态射分别定义 \(R^nF(f)\) 或 \(L_nF(f)\)。
\end{definition}

由\cref{b4:prop:comparison-homotopy-unique,b4:cor:resolutions-homotopy-equivalent}，提升同一个态射的比较映射同伦，因而诱导相同的同调态射；恒等和复合也由比较映射的恒等和复合诱导。于是上述定义给出函子，更换消解只引起满足传递律的自然同构。它们还是加性函子：若 \(f,g:A\to A'\) 分别由比较映射 \(\widetilde f,\widetilde g\) 提升，则 \(\widetilde f+\widetilde g\) 提升 \(f+g\)。由于 \(F\) 加性，取同调后有
\[
R^nF(f+g)=R^nF(f)+R^nF(g),\qquad
L_nF(f+g)=L_nF(f)+L_nF(g).
\]
若把投射消解写成 \(P^{-n}=P_n\)，则 \(L_nF(A)=H^{-n}(F(P^\bullet))\)。下标 \(n\) 是非负同调次数，不是把右导出函子也改放到负次数。

\begin{proposition}\label{b4:prop:derived-degree-zero}
设 \(\mathcal A,\mathcal B\) 是交换范畴。若 \(\mathcal A\) 有足够多内射对象且 \(F:\mathcal A\to\mathcal B\) 加性、左正合，则有自然同构 \(R^0F\cong F\)，且内射对象 \(I\) 满足 \(R^nF(I)=0\) 对 \(n>0\)。若有足够多投射对象且 \(F\) 加性、右正合，则 \(L_0F\cong F\)，且投射对象的正次数左导出函子消失。
\end{proposition}
\begin{proof}
由 \(0\to A\to I^0\to I^1\) 正合及左正合性，\(F(A)\) 就是 \(F(I^0)\to F(I^1)\) 的核，即零次上同调；构造与态射相容。内射对象可取自身集中于零次的内射消解，故正次消失。投射情形用 \(P_1\to P_0\to A\to0\) 的余核及零次链消解。
\end{proof}

\cref{b4:tab:derived-conventions}集中列出两种导出的消解与次数约定，其中左、右正合性保证零次导出与原函子相同。表中的连接态射对应短正合列 \(0\to A\to B\to C\to0\)：右导出的连接升一次，适用于 \(n\ge0\)；左导出的连接降一次，适用于 \(n\ge1\)。模范畴中的 Ext 属于右导出，Tor 属于左导出；Ext 的第一变量反变，还需在相应长正合列中反转对象的顺序。

\begin{table}[htbp]
\centering
\caption{左、右导出函子的约定}
\label{b4:tab:derived-conventions}
\begin{tabular}{@{}lllll@{}}
\toprule
类型 & \(F\) 的条件 & 所选消解 & 记号 & 连接态射 \\
\midrule
右导出 & 加性、左正合 & 内射上链消解 & \(R^nF\)
& \(R^nF(C)\to R^{n+1}F(A)\) \\
左导出 & 加性、右正合 & 投射链消解 & \(L_nF\)
& \(L_nF(C)\to L_{n-1}F(A)\) \\
\bottomrule
\end{tabular}
\end{table}


\cref{b4:prop:horseshoe-functor-sequence}给出短正合列的导出长正合列。下一小节把这项性质与消解的具体选取分开表述。

\subsection{\texorpdfstring{\(\delta\)-函子与泛性质}{\AlgebraPdfDelta-函子与泛性质}}
\label{b4:subsec:0801:02}

\begin{definition}\label{b4:def:delta-functor}
设 \(\mathcal A,\mathcal B\) 是交换范畴。一族加性函子 \(T^n:\mathcal A\to\mathcal B\)，\(n\ge0\)，若对每个短正合列 \(0\to A\to B\to C\to0\) 配有态射 \(\delta^n:T^n(C)\to T^{n+1}(A)\)，使
\[
0\longrightarrow T^0(A)\longrightarrow T^0(B)\longrightarrow T^0(C)
\xrightarrow{\delta^0}T^1(A)\longrightarrow\cdots
\]
正合，且连接态射对短正合列之间的态射自然，就称 \(T^\bullet\) 连同这些连接态射为\term[delta-hanzi]{上同调 \(\delta\)-函子}[cohomological delta functor]。若 \(T^\bullet,U^\bullet:\mathcal A\to\mathcal B\) 均为此类函子，且逐次自然变换 \(a^n:T^n\to U^n\) 与每个连接态射交换，则称 \(a^\bullet\) 为它们之间的态射。
\end{definition}

对短正合列 \(0\to A\to B\to C\to0\)，\(\delta\)-函子态射 \(a^\bullet:T^\bullet\to U^\bullet\) 所要求的相容性就是下列方块交换：
\[
\begin{tikzcd}[column sep=large,row sep=large]
T^n(C)\arrow[r,"\delta_T^n"]\arrow[d,"a_C^n"']&T^{n+1}(A)\arrow[d,"a_A^{n+1}"]\\
U^n(C)\arrow[r,"\delta_U^n"']&U^{n+1}(A)
\end{tikzcd}
\]
同调版本由 \(T_n\) 和连接态射 \(T_n(C)\to T_{n-1}(A)\) 组成，长正合列在右端以 \(T_0(C)\to0\) 结束。右导出函子组成上同调 \(\delta\)-函子，左导出函子组成同调 \(\delta\)-函子。

\begin{definition}\label{b4:def:universal-delta-functor}
设 \(\mathcal A,\mathcal B\) 为交换范畴，\(T^\bullet:\mathcal A\to\mathcal B\) 为上同调 \(\delta\)-函子。若对任意上同调 \(\delta\)-函子 \(U^\bullet:\mathcal A\to\mathcal B\)，每个自然变换 \(T^0\to U^0\) 都唯一延拓为 \(\delta\)-函子态射 \(T^\bullet\to U^\bullet\)，则称 \(T^\bullet\) 为\term[fan-delta-hanzi]{泛上同调 \(\delta\)-函子}[universal cohomological delta functor]。对同调 \(\delta\)-函子 \(T_\bullet:\mathcal A\to\mathcal B\)，若对任意同调 \(\delta\)-函子 \(U_\bullet:\mathcal A\to\mathcal B\)，每个 \(U_0\to T_0\) 都唯一延拓为 \(\delta\)-函子态射 \(U_\bullet\to T_\bullet\)，则称 \(T_\bullet\) 具有同调泛性。
\end{definition}

这里方向不能省略。右导出函子从自身的零次项出发控制到其他理论的映射；左导出函子接收其他理论到自身零次项的映射。它们分别由余核式和核式的归纳构造产生。相同零次函子的两个泛 \(\delta\)-函子必自然同构：双向延拓零次恒等，再用唯一性证明两个复合都是恒等。

\subsection{\texorpdfstring{消去性与泛 \(\delta\)-函子}{消去性与泛 \AlgebraPdfDelta-函子}}
\label{b4:subsec:0801:03}

\begin{definition}\label{b4:def:effaceable}
设 \(T:\mathcal A\to\mathcal B\) 是交换范畴间的加性函子。若每个对象 \(A\) 都有单态射 \(j:A\to E\) 使 \(T(j)=0\)，则称 \(T\) \term[xiaoqvxing]{可消去}[effaceable]。若每个对象 \(A\) 都有满态射 \(P\to A\) 使其在 \(T\) 下为零，则称 \(T\) 可对偶消去。
\end{definition}

消去的是一个诱导态射，不一定要求 \(T(E)=0\)。导出函子的消去性更强：嵌入内射对象便使其正次取值本身为零。

\begin{theorem}[消去性与泛性]\label{b4:thm:derived-universal}
设 \(\mathcal A,\mathcal B\) 为交换范畴，\(T^\bullet:\mathcal A\to\mathcal B\) 为上同调 \(\delta\)-函子。若每个 \(T^n\)、\(n>0\)，均可消去，则它是泛上同调 \(\delta\)-函子。对偶地，正次均可对偶消去的同调 \(\delta\)-函子具有同调泛性。特别地，对加性函子 \(F:\mathcal A\to\mathcal B\)，若 \(\mathcal A\) 有足够多内射对象且 \(F\) 左正合，则 \(R^\bullet F\) 具有上同调泛性；若 \(\mathcal A\) 有足够多投射对象且 \(F\) 右正合，则 \(L_\bullet F\) 具有同调泛性。
\end{theorem}
\begin{proof}
给定 \(a^0:T^0\to U^0\)。归纳假定次数小于 \(n\) 的自然变换已定义，并与已有连接态射相容。对 \(A\) 取单态射 \(j:A\to E\)，使 \(T^n(j)=0\)，记 \(Q=E/A\)。正合性给出满态射
\[
\delta_T:T^{n-1}(Q)\longrightarrow T^n(A),
\qquad \ker\delta_T=\operatorname{im}\bigl(T^{n-1}(E)\to T^{n-1}(Q)\bigr).
\]
复合 \(\delta_Ua_Q^{n-1}\) 在这个核上为零：由低次自然性，其在 \(T^{n-1}(E)\) 上的值等于先用 \(a_E^{n-1}\)，再沿 \(U\) 的相邻两个态射复合，而后者为零。因此它唯一下降为
\[
a_A^n:T^n(A)\longrightarrow U^n(A),
\qquad a_A^n\delta_T=\delta_Ua_Q^{n-1}.
\]
这里 \(\delta_T\) 是满态射，所以等式唯一决定 \(a_A^n\)；任何与连接态射相容的延拓都必须取此值。

检验选择独立性。若另选 \(j':A\to E'\)，则 \(j''=(j,j'):A\to E\oplus E'\) 仍单，且被 \(T^n\) 零化。记 \(Q''=\operatorname{coker}j''\)。向 \(E\) 的投影 \(\pi_E\) 诱导下图，其中两行短正合：
\[
\begin{tikzcd}[column sep=small,row sep=large]
0\arrow[r]&A\arrow[r,"j''"]\arrow[d,"1_A"']
&E\oplus E'\arrow[r]\arrow[d,"\pi_E"]
&Q''\arrow[r]\arrow[d,"\bar\pi_E"]&0\\
0\arrow[r]&A\arrow[r,"j"']&E\arrow[r]&Q\arrow[r]&0.
\end{tikzcd}
\]
向 \(E'\) 的投影给出同样的比较图。用下标 \(j,j''\) 区分两种选择所得的候选态射。低次自然性及连接态射的自然性给出
\[
a_{A,j}^n\delta_T''=\delta_U''a_{Q''}^{n-1}
=a_{A,j''}^n\delta_T''.
\]
上行的 \(\delta_T''\) 为满态射，故 \(a_{A,j}^n=a_{A,j''}^n\)。对另一投影重复此论证，得到 \(a_{A,j'}^n=a_{A,j''}^n\)，从而构造不依赖选择。

检验对任意 \(f:A\to B\) 的自然性。选分别消去 \(T^n(A),T^n(B)\) 的 \(j_A:A\to E_A\)、\(j_B:B\to E_B\)。用 \(j_f=(j_A,j_B\circ f):A\to E_A\oplus E_B\) 作为 \(A\) 的选择，记 \(Q_f=\operatorname{coker}j_f\)、\(Q_B=\operatorname{coker}j_B\)。向 \(E_B\) 的投影诱导
\[
\begin{tikzcd}[column sep=small,row sep=large]
0\arrow[r]&A\arrow[r,"j_f"]\arrow[d,"f"']
&E_A\oplus E_B\arrow[r]\arrow[d,"\pi_B"]
&Q_f\arrow[r]\arrow[d,"\bar f"]&0\\
0\arrow[r]&B\arrow[r,"j_B"']&E_B\arrow[r]&Q_B\arrow[r]&0.
\end{tikzcd}
\]
已证的选择独立性允许用上行定义 \(a_A^n\)。将低次自然性与连接态射的自然性用于此图，两种复合满足
\[
\bigl(U^n(f)a_A^n\bigr)\delta_{T,f}
=\delta_{U,B}a_{Q_B}^{n-1}T^{n-1}(\bar f)
=\bigl(a_B^nT^n(f)\bigr)\delta_{T,f}.
\]
消去满态射 \(\delta_{T,f}\)，即得 \(U^n(f)a_A^n=a_B^nT^n(f)\)。

最后检验与任意 \(0\to A\xrightarrow{i}B\to C\to0\) 的连接态射相容。选择 \(j:B\to E\) 消去 \(T^n(B)\)，则 \(ji:A\to E\) 消去 \(T^n(A)\)。原短正合列映入 \(0\to A\to E\to E/A\to0\)，其中左端为恒等。定义式已保证后一列的相容；把它与 \(C\to E/A\) 的低次自然性复合，就得到原列的相容。归纳完成存在、自然性及唯一性。

把论证同时用于 \(\mathcal A^{\mathrm{op}}\) 与 \(\mathcal B^{\mathrm{op}}\)，单态射变满态射，余核下降变核分解，自然变换的方向也反转，遂得同调版本。最后，\cref{b4:prop:derived-degree-zero}使内射嵌入、投射满射消去所有正次导出函子，故可应用已证判据。
\end{proof}

这一泛性质确定的不仅是各个同调对象，也包括所有态射和连接态射。以后从另一种消解或另一种构造得到同一理论时，可以用零次的自然同构来识别整套长正合列。

% !TeX root = ../../Book4.tex
% 纲要标题：### 9.2 Ext
% 纲要标签：b4:sec:0802
% 纲要位置：计划纲要.md，第 3408 行。

\section{Ext}
\label{b4:sec:0802}

Hom 的两个变量具有相反的箭头方向，分别用投射或内射消解得到的计算必须比较。平衡性使两种模型计算同一组 Ext，并保留两变量的连接态射。


% 纲要内容（仅作写作计划，尚非教材正文）：
%
% * Hom 的右导出函子
% * 反变与协变变量
% * Ext 长正合列及具体计算
%

\subsection{Hom 的右导出函子}
\label{b4:subsec:0802:01}

\begin{definition}\label{b4:def:ext}
设 \(R\) 是含幺环，\(M,N\) 是幺性左 \(R\) 模。对 \(N\) 取内射消解 \(N\to I^\bullet\)，定义交换群
\[
\operatorname{Ext}_R^n(M,N)\coloneqq
H^n\bigl(\operatorname{Hom}_R(M,I^\bullet)\bigr),\qquad n\ge0.
\]
这就是左正合函子 \(\operatorname{Hom}_R(M,-)\) 的右导出函子。
\end{definition}
\symidx{\(\operatorname{Ext}_R^n(M,N)\)：Hom 函子的右导出函子值}

内射消解的开头是 \(0\to N\xrightarrow{\eta}I^0\xrightarrow{\mathrm{d}_I^0}I^1\)。因此
\[
\operatorname{Ext}_R^0(M,N)
=\ker\bigl(\operatorname{Hom}_R(M,I^0)
\xrightarrow{\mathrm{d}_I^0\circ-}\operatorname{Hom}_R(M,I^1)\bigr)
\cong\operatorname{Hom}_R(M,N).
\]
这里同构由 \(f\mapsto\eta f\) 给出：映入 \(I^0\) 后被微分零化的态射，恰好唯一经过核 \(\eta:N\hookrightarrow I^0\)。

\ProseParagraphBreak
在一般非交换环上，Ext 首先是交换群，不自动是左 \(R\) 模。即使在零次 Hom 中，试图定义 \((rf)(m)=r f(m)\) 也可能破坏 \(R\) 线性，因为
\[
(rf)(sm)=rs f(m),\qquad s(rf)(m)=sr f(m),
\]
两者未必相等。若 \(R\) 交换，这种标量乘法保持线性，并与 Hom 复形的所有微分相容，便得到自然的 \(R\) 模结构。下标的环也不能略去，例如 \(\operatorname{Ext}_{\mathbb Z}^1(\mathbb Z/2\mathbb Z,\mathbb Z/2\mathbb Z)\) 与 \(\operatorname{Ext}_{\mathbb F_2}^1(\mathbb F_2,\mathbb F_2)\) 不同。

\ProseParagraphBreak
固定 \(N\) 时，\(\operatorname{Hom}_R(-,N)\) 是从左模反范畴到交换群的左正合函子。左模的投射对象正是反范畴的内射对象；原来的投射消解
\[
\cdots\longrightarrow P_1\longrightarrow P_0\longrightarrow M\longrightarrow0
\]
在反范畴中成为 \(0\to M\to P_0\to P_1\to\cdots\)，即其中的一条内射消解。施加这个反变 Hom 后，\(\operatorname{Hom}_R(P_n,N)\) 的次数从 \(n\) 向 \(n+1\) 增加，所以它给出另一套右导出函子。两套计算是否一致，需要证明，不能由名字相同得出。

\subsection{反变与协变变量}
\label{b4:subsec:0802:02}

为了方便表示上同调类，以下把投射消解 \(P_\bullet\to M\) 送入 Hom 时，采用无符号的微分 \(\delta^n(f)=f\mathrm{d}_{P,n+1}\)。这与\cref{b4:def:hom-complex}的约定可以相互转换：若把 \(P_n\) 放在上链次数 \(-n\)，把 \(N\) 放在零次，则内部 Hom 的微分是 \(\delta_{\mathrm{Hom}}^n(f)=(-1)^{n+1}\cdot f\mathrm{d}_{P,n+1}\)。令 \(c_n=(-1)^{n(n+1)/2}\)，在次数 \(n\) 作重标度 \(\rho_n(f)=c_nf\)，则
\[
c_{n+1}=(-1)^{n+1}c_n,
\qquad
\rho_{n+1}\delta_{\mathrm{Hom}}^n(f)
=c_{n+1}(-1)^{n+1}\cdot f\mathrm{d}_{P,n+1}
=\delta^n\rho_n(f).
\]
因此 \(\rho\) 是从内部 Hom 复形到无符号复形的同构。例如 \(c_0,c_1,c_2,c_3\) 依次为 \(1,-1,-1,1\)，恰好吸收微分在相邻次数的符号。本节以后均采用无符号的计算约定。

\begin{theorem}[Ext 的平衡性]\label{b4:thm:ext-balanced}
设 \(R\) 是含幺环，\(M,N\) 是幺性左 \(R\) 模，\(P_\bullet\to M\) 为投射消解，\(N\to I^\bullet\) 为内射消解。对整数 \(n\ge0\) 取
\[
\operatorname{Hom}_R(P_\bullet,N)^n=\operatorname{Hom}_R(P_n,N),
\qquad \delta^n(f)=f\,\mathrm{d}_{P,n+1}.
\]
则有自然同构
\[
H^n\bigl(\operatorname{Hom}_R(P_\bullet,N)\bigr)
\xrightarrow[\sim]{\beta_n}\operatorname{Ext}_R^n(M,N).
\]
它对 \(M\) 反变，对 \(N\) 协变，并与两变量的连接态射相容。
\end{theorem}
\begin{proof}
两种计算将通过同一个双复形比较。令 \(D^{p,q}=\operatorname{Hom}_R(P_p,I^q)\)，其中 \(p,q\ge0\)：固定 \(p\) 沿内射消解计算，固定 \(q\) 则沿投射消解计算。反变 Hom 把 \(P_{p+1}\to P_p\) 变为 \(D^{p,q}\to D^{p+1,q}\)，所以这里的两个指标都按上链次数增加。定义
\[
\mathrm{d}_h(f)=f\mathrm{d}_{P,p+1},\qquad
\mathrm{d}_v(f)=(-1)^p\cdot\mathrm{d}_I^qf.
\]
两次水平或竖直微分为零；一水平一竖直的两种复合中，符号分别为 \((-1)^p\) 和 \((-1)^{p+1}\)，故 \(\mathrm{d}_h\mathrm{d}_v+\mathrm{d}_v\mathrm{d}_h=0\)。总次数为 \(p+q\)。对 \(n\ge0\)，\(p+q=n\) 只有 \(n+1\) 对非负整数解；对 \(n<0\) 则没有解，故总复形逐对角有限。

固定 \(p\)，\(P_p\) 的投射性使 \(\operatorname{Hom}_R(P_p,-)\) 正合。把内射消解的开头送入它，得到
\[
0\longrightarrow\operatorname{Hom}_R(P_p,N)
\longrightarrow D^{p,0}\longrightarrow D^{p,1},
\]
故纵列零次上同调为 \(\ker(D^{p,0}\to D^{p,1})\cong\operatorname{Hom}_R(P_p,N)\)，正次数上同调为零。固定 \(q\)，\(I^q\) 的内射性同样给出横行零次上同调 \(\ker(D^{0,q}\to D^{1,q})\cong\operatorname{Hom}_R(M,I^q)\)，正次数上同调为零。因此需要比较的两个 Hom 复形分别由下边缘的纵核和左边缘的横核组成，而不是双复形中整条底行或整条左列。

记 \(\eta:N\to I^0\)、\(\varepsilon:P_0\to M\) 为内射、投射消解的增广。它们分别给出复形映射
\[
\begin{aligned}
\operatorname{Hom}_R(P_\bullet,N)&\longrightarrow\operatorname{Tot}D,
&f&\longmapsto\eta f\in D^{p,0},\\
\operatorname{Hom}_R(M,I^\bullet)&\longrightarrow\operatorname{Tot}D,
&g&\longmapsto g\varepsilon\in D^{0,q}.
\end{aligned}
\]
第一条由 \(N\to I^0\) 后合成而来，\(\mathrm{d}_I^0\eta=0\) 使其像的竖直微分为零；第二条由 \(P_0\to M\) 预合成而来，\(\varepsilon\mathrm{d}_{P,1}=0\) 使其像的水平微分为零。其余微分正是边缘复形自身的微分，所以两条映射与总微分相容。将边缘复形分别放在 \(q=0\) 和 \(p=0\)，刚才的列、行上同调计算说明这两个双复形映射分别逐列、逐行拟同构。\cref{b4:cor:double-complex-comparison}的有限对角条件已经满足，故两条总化映射都是拟同构。

将 \(P_p\) 放在上链次数 \(-p\)，总复形到内部 Hom 复形的同构在分量 \(D^{p,q}\) 上为
\[
 S^{p,q}(f)=c_p(-1)^{pq}f.
\]
沿水平方向，其系数之比为 \((-1)^{p+q+1}\)；沿竖直方向，其系数之比为 \((-1)^p\)，恰好把上面的总微分变成内部 Hom 微分。左边缘 \(p=0\) 的系数为 \(1\)，下边缘 \((p,q)=(n,0)\) 的系数为 \(c_n\)。因此 \(\beta_n\) 将无符号余圈 \(b:P_n\to N\) 对应于内部 Hom 余圈 \(c_nb\)，这一系数来自两条边缘映射的比较。

投射、内射比较映射使上述两条映射对 \(M,N\) 自然；不同选择同伦，故同调态射确定。还需比较连接态射：两种模型不仅应得到同构的交换群，也应得到同一组长正合列，否则不能无歧义地使用 Ext 的连接态射。分别在两个变量中核对这一点。给定第一变量的短正合列 \(0\to A\to B\to C\to0\)，取投射马蹄消解
\[
0\longrightarrow P_{A,\bullet}\longrightarrow P_{B,\bullet}
\longrightarrow P_{C,\bullet}\longrightarrow0,
\]
它在每个次数分裂。固定 \(N\) 的内射消解，对 \(X=A,B,C\) 令 \(T_X=\operatorname{Tot}\operatorname{Hom}_R(P_{X,\bullet},I^\bullet)\)。Hom 的第一变量反变，把 \(A,B,C\) 的次序反转为 \(C,B,A\)，因此得到的短正合列方向是
\[
0\longrightarrow T_C\longrightarrow T_B\longrightarrow T_A\longrightarrow0.
\]
两侧边缘复形也分别组成 \(C,B,A\) 顺序的短正合列：投射计算一侧用马蹄消解的逐次分裂性，内射计算一侧用各 \(I^q\) 的内射性。边缘到总复形的增广映射给出短正合列之间的交换图。由\cref{b4:thm:cohomology-long-exact}的自然性，两种计算中的连接态射
\[
\operatorname{Ext}_R^n(A,N)\longrightarrow\operatorname{Ext}_R^{n+1}(C,N)
\]
在平衡同构下相互对应。

给定第二变量的短正合列 \(0\to N'\to N\to N''\to0\)，取逐次分裂的内射马蹄消解 \(0\to I_{N'}^\bullet\to I_N^\bullet\to I_{N''}^\bullet\to0\)。固定 \(M\) 的投射消解，令 \(T_V=\operatorname{Tot}\operatorname{Hom}_R(P_\bullet,I_V^\bullet)\)，其中 \(V=N',N,N''\)。这次 Hom 保持箭头方向，得到
\[
0\longrightarrow T_{N'}\longrightarrow T_N\longrightarrow T_{N''}\longrightarrow0.
\]
投射计算一侧的边缘短正合列使用各 \(P_p\) 的投射性，内射计算一侧则使用马蹄消解的逐次分裂性。增广仍给出短正合列之间的交换图，故连接态射
\[
\operatorname{Ext}_R^n(M,N'')\longrightarrow\operatorname{Ext}_R^{n+1}(M,N')
\]
也与平衡同构相容。
\end{proof}

给定 \(u:M'\to M\)，其投射消解提升经预合成给出 \(u^*:\operatorname{Ext}_R^n(M,N)\to\operatorname{Ext}_R^n(M',N)\)；给定 \(v:N\to N'\)，后合成给出 \(v_*\)。预合成与后合成交换，所以得到双函子
\[
\operatorname{Ext}_R^n:
(R\text{-}\mathbf{Mod})^{\mathrm{op}}\times R\text{-}\mathbf{Mod}
\longrightarrow\mathbf{Ab}.
\]
它同时记录第一变量的反变作用与第二变量的协变作用。

\subsection{Ext 长正合列及具体计算}
\label{b4:subsec:0802:03}

\begin{theorem}\label{b4:thm:ext-long-exact}
设 \(R\) 是含幺环。幺性左 \(R\) 模短正合列 \(0\to A\to B\to C\to0\) 对任意幺性左模 \(N\) 给出自然长正合列
\[
\begin{aligned}
0&\longrightarrow\operatorname{Hom}_R(C,N)\longrightarrow\operatorname{Hom}_R(B,N)
\longrightarrow\operatorname{Hom}_R(A,N)\\
&\longrightarrow\operatorname{Ext}_R^1(C,N)\longrightarrow\operatorname{Ext}_R^1(B,N)
\longrightarrow\operatorname{Ext}_R^1(A,N)\longrightarrow\cdots.
\end{aligned}
\]
其一般中间段为
\[
\begin{aligned}
\cdots\longrightarrow\operatorname{Ext}_R^n(C,N)
&\longrightarrow\operatorname{Ext}_R^n(B,N)
\longrightarrow\operatorname{Ext}_R^n(A,N)\\
&\longrightarrow\operatorname{Ext}_R^{n+1}(C,N)\longrightarrow\cdots
\qquad(n\ge1).
\end{aligned}
\]
对任意幺性左模 \(M\)，幺性左模短正合列 \(0\to N'\to N\to N''\to0\) 给出自然长正合列
\[
\begin{aligned}
0&\longrightarrow\operatorname{Hom}_R(M,N')\longrightarrow\operatorname{Hom}_R(M,N)
\longrightarrow\operatorname{Hom}_R(M,N'')\\
&\longrightarrow\operatorname{Ext}_R^1(M,N')\longrightarrow\operatorname{Ext}_R^1(M,N)
\longrightarrow\operatorname{Ext}_R^1(M,N'')\longrightarrow\cdots.
\end{aligned}
\]
其一般中间段为
\[
\begin{aligned}
\cdots\longrightarrow\operatorname{Ext}_R^n(M,N')
&\longrightarrow\operatorname{Ext}_R^n(M,N)
\longrightarrow\operatorname{Ext}_R^n(M,N'')\\
&\longrightarrow\operatorname{Ext}_R^{n+1}(M,N')\longrightarrow\cdots
\qquad(n\ge1).
\end{aligned}
\]
\end{theorem}
\begin{proof}
第一列取 \(N\) 的内射消解，逐次施加反变正合函子 \(\operatorname{Hom}_R(-,I^q)\)，得到 \(0\to\operatorname{Hom}(C,I)\to\operatorname{Hom}(B,I)\to\operatorname{Hom}(A,I)\to0\)，再取同调长正合列。第二列取 \(M\) 的投射消解，对给定序列逐次施加正合函子 \(\operatorname{Hom}_R(P_p,-)\)，取长正合列，再用\cref{b4:thm:ext-balanced}识别各项。增广在零次给出通常 Hom 映射，连接态射与选择独立且自然。
\end{proof}

\begin{example}\label{b4:exm:ext-cyclic}
设 \(R\) 是主理想整环，\(0\ne a\in R\)，\(N\) 是幺性 \(R\) 模。\(R/(a)\) 的两项自由消解为 \(0\to R\xrightarrow{\times a}R\to R/(a)\to0\)。送入 Hom 后得到
\[
0\longrightarrow N\xrightarrow{\times a}N\longrightarrow0,
\]
两个 \(N\) 分别位于上链次数零、一；零次上同调是乘 \(a\) 的核，一次上同调是其余核。因此
\[
\operatorname{Ext}_R^0(R/(a),N)=\{\,n\in N\mid an=0\,\},
\qquad\operatorname{Ext}_R^1(R/(a),N)=N/aN,
\]
二次及以上为零。特别地，\(\operatorname{Ext}_{\mathbb Z}^1(\mathbb Z/m\mathbb Z,\mathbb Z/n\mathbb Z)\cong\mathbb Z/\gcd(m,n)\mathbb Z\)，其中 \(m,n>0\)。若改取 \(N=\mathbb Q\)，从这个具体复形可见乘 \(m\) 为同构，所以零次及正次均为零。另一方面，\(\mathbb Q\) 的内射性也保证正次数 Ext 消失，与计算相符。
\end{example}

设 \(k\) 为域，在 \(R=k[\epsilon]/(\epsilon^2)\) 上，对\secref{b4:sec:0701}给出的 \(k=R/(\epsilon)\) 的周期消解施加 \(\operatorname{Hom}_R(-,k)\)，所有微分为零。因此 \(\operatorname{Ext}_R^n(k,k)\cong k\) 对每个 \(n\ge0\) 成立。相同的底层向量空间，在底环为 \(k\) 时正次 Ext 全为零；记录底环是计算的一部分。

% !TeX root = ../../Book4.tex
% 纲要标题：### 9.3 Tor
% 纲要标签：b4:sec:0803
% 纲要位置：计划纲要.md，第 3414 行。

\section{Tor}
\label{b4:sec:0803}

张量积的两个模必须取相反的侧别，任一变量的投射消解都应得到同一组 Tor。总复形的比较同时解释平衡性和短正合列中的连接映射。


% 纲要内容（仅作写作计划，尚非教材正文）：
%
% * 张量积的左导出函子
% * 平衡性
% * Tor 长正合列及具体计算
%

\subsection{张量积的左导出函子}
\label{b4:subsec:0803:01}

\begin{definition}\label{b4:def:tor}
设 \(R\) 是含幺环，\(M\) 是幺性右 \(R\) 模，\(N\) 是幺性左 \(R\) 模。对 \(M\) 取右模投射消解 \(P_\bullet\to M\)，定义交换群
\[
\operatorname{Tor}_n^R(M,N)\coloneqq H_n(P_\bullet\otimes_R N),\qquad n\ge0.
\]
这就是右正合函子 \(-\otimes_R N\) 的左导出函子。
\end{definition}
\symidx{\(\operatorname{Tor}_n^R(M,N)\)：张量积函子的左导出函子值}

投射消解在零次给出 \(P_1\to P_0\to M\to0\)。张量的右正合性使
\[
P_1\otimes_RN\longrightarrow P_0\otimes_RN
\longrightarrow M\otimes_RN\longrightarrow0
\]
正合，因此
\[
\operatorname{Tor}_0^R(M,N)
=\operatorname{coker}(P_1\otimes_RN\to P_0\otimes_RN)
\cong M\otimes_RN.
\]
这一零次识别只用右正合性，不要求 \(N\) 平坦。

\ProseParagraphBreak
左右侧别不可交换省略。若 \(P_\bullet\to M\) 和 \(Q_\bullet\to N\) 分别为右模、左模投射消解，下一小节的平衡性给出
\[
H_n(P_\bullet\otimes_RN)\cong H_n(M\otimes_RQ_\bullet).
\]
它没有交换 \(M,N\)，只改变在哪个变量取消解。只有在交换环上，左右模有自然识别，才能进一步比较 \(\operatorname{Tor}_n^R(M,N)\) 与 \(\operatorname{Tor}_n^R(N,M)\)，得到对称性。

\subsection{平衡性}
\label{b4:subsec:0803:02}

\begin{theorem}[Tor 的平衡性]\label{b4:thm:tor-balanced}
设 \(R\) 是含幺环，\(M\) 是幺性右 \(R\) 模，\(N\) 是幺性左 \(R\) 模。若 \(Q_\bullet\to N\) 是左 \(R\) 模投射消解，则对每个整数 \(n\ge0\)，有对两个变量自然的同构
\[
\operatorname{Tor}_n^R(M,N)\cong H_n(M\otimes_RQ_\bullet).
\]
因而可以在任意一侧取投射消解计算 Tor。
\end{theorem}
\begin{proof}
与 Ext 的平衡性一样，取两份消解组成双复形，再比较其两条边缘。这里使行、列正合的性质是投射模的平坦性。取右模投射消解 \(P_\bullet\to M\)，令 \(D_{p,q}=P_p\otimes_RQ_q\)，\(p,q\ge0\)，以
\[
\mathrm{d}_h=\mathrm{d}_P\otimes1,\qquad \mathrm{d}_v=(-1)^p1\otimes \mathrm{d}_Q
\]
为微分。它们反交换，所以总链微分平方为零。第二卷定理7.3.4证明了投射左模平坦。投射右 \(R\) 模正是投射左 \(R^{\mathrm{op}}\) 模，将该定理用于 \(R^{\mathrm{op}}\)，并经 \(X\otimes_{R^{\mathrm{op}}}P\cong P\otimes_RX\) 识别相应的张量函子，便知 \(P\otimes_R-\) 正合，即投射右模平坦。固定 \(p\) 时，\(P_p\) 是投射右 \(R\) 模，因而 \(P_p\otimes_R-\) 在左模范畴上正合；竖列只在零次有同调 \(P_p\otimes_RN\)。固定 \(q\) 时，\(Q_q\) 是投射左 \(R\) 模，因而 \(-\otimes_RQ_q\) 在右模范畴上正合；横行只在零次有同调 \(M\otimes_RQ_q\)。

记两份消解的增广为 \(\varepsilon_P:P_0\to M\)、\(\varepsilon_Q:Q_0\to N\)。第一条总化映射在 \(q=0\) 上取 \(1\otimes\varepsilon_Q\)，在 \(q>0\) 上取零；第二条在 \(p=0\) 上取 \(\varepsilon_P\otimes1\)，在 \(p>0\) 上取零。由增广零化对应的一次微分，它们与总链微分交换，给出
\[
\operatorname{Tot}D\longrightarrow P_\bullet\otimes_RN,
\qquad
\operatorname{Tot}D\longrightarrow M\otimes_RQ_\bullet.
\]
把两条边缘分别放在 \(q=0\) 和 \(p=0\)，刚才的纵列、横行同调计算说明这两条映射分别逐列、逐行拟同构。对于 \(n\ge0\)，总次数对角线 \(p+q=n\) 只有 \((0,n),(1,n-1),\ldots,(n,0)\) 共 \(n+1\) 个位置；\(n<0\) 时没有位置。因此\cref{b4:cor:double-complex-comparison}的链版本适用，两条总化映射均为拟同构。取同调并组合两同构即得结论。比较映射与增广交换，比较选择同伦，故所得同构对两变量自然。

还需要使两种计算给出相同的连接态射。对于第一变量的短正合列，取逐次分裂的右模马蹄消解，并固定第二变量的左模投射消解。三个总复形及边缘 \(P_\bullet\otimes_RN\) 的短正合列使用马蹄的逐次分裂性；另一边缘 \(M\otimes_RQ_\bullet\) 的短正合列则使用各 \(Q_q\) 平坦。对于第二变量的短正合列，改用左模马蹄消解并固定 \(P_\bullet\)，两条边缘的正合性分别来自各 \(P_p\) 平坦和马蹄的逐次分裂性。两种情况下，增广映射都给出总复形短正合列到边缘短正合列的交换图。由同调长正合列的自然性，平衡同构便与两变量的连接态射相容。
\end{proof}

若 \(R\) 交换，对左模定义右作用 \(n\cdot r\coloneqq rn\)，对右模定义左作用 \(r\cdot m\coloneqq mr\)，便可自然识别左右模；同样识别消解中的各项。这时交换两因子的映射才有如下张量积类型，并且满足平衡关系：
\[
\tau_{p,q}:P_p\otimes_RQ_q\longrightarrow Q_q\otimes_RP_p,
\qquad x\otimes y\longmapsto(-1)^{pq}y\otimes x.
\]
这些 \(\tau_{p,q}\) 合为保持总次数的映射 \(\tau\)。在交换后以 \(Q_q\) 为第一因子的总复形中，微分为 \(\mathrm{d}'=\mathrm{d}_Q\otimes1+(-1)^q1\otimes\mathrm{d}_P\)。对 \(x\in P_p\)、\(y\in Q_q\)，有
\[
\begin{aligned}
\mathrm{d}'\tau(x\otimes y)
&=(-1)^{pq}\cdot\mathrm{d}_Qy\otimes x
 +(-1)^{pq+q}y\otimes\mathrm{d}_Px,\\
\tau\mathrm{d}(x\otimes y)
&=(-1)^{(p-1)q}y\otimes\mathrm{d}_Px
 +(-1)^{p+p(q-1)}\cdot\mathrm{d}_Qy\otimes x.
\end{aligned}
\]
其中 \(p+p(q-1)=pq\)，而 \((p-1)q\) 与 \(pq+q\) 相差 \(2q\)，故两个分量的系数分别相同，得到 \(\mathrm{d}'\tau=\tau\mathrm{d}\)。因此 \(\tau\) 给出总复形同构，取同调并用平衡性，得到对两变量自然的同构 \(\operatorname{Tor}_n^R(M,N)\cong\operatorname{Tor}_n^R(N,M)\)。一般非交换环上，仅凭给定的左右模结构，不能作上述自然识别或定义这样的交换映射；平衡性仍只表示可在任一变量取消解。

\subsection{Tor 长正合列及具体计算}
\label{b4:subsec:0803:03}

\begin{theorem}\label{b4:thm:tor-long-exact}
设 \(R\) 是含幺环，\(M\) 是幺性右 \(R\) 模，\(N\) 是幺性左 \(R\) 模。幺性右模短正合列 \(0\to M'\to M\to M''\to0\) 给出
\[
\begin{aligned}
\cdots&\longrightarrow\operatorname{Tor}_1^R(M,N)
\longrightarrow\operatorname{Tor}_1^R(M'',N)
\longrightarrow M'\otimes_RN\\
&\longrightarrow M\otimes_RN\longrightarrow M''\otimes_RN\longrightarrow0,
\end{aligned}
\]
其一般中间段为
\[
\begin{aligned}
\cdots\longrightarrow\operatorname{Tor}_n^R(M',N)
&\longrightarrow\operatorname{Tor}_n^R(M,N)
\longrightarrow\operatorname{Tor}_n^R(M'',N)\\
&\xrightarrow{\delta}\operatorname{Tor}_{n-1}^R(M',N)
\longrightarrow\cdots\qquad(n\ge1).
\end{aligned}
\]
幺性左 \(R\) 模短正合列 \(0\to N'\to N\to N''\to0\) 同样给出自然长正合列，其一般中间段为
\[
\begin{aligned}
\cdots\longrightarrow\operatorname{Tor}_n^R(M,N')
&\longrightarrow\operatorname{Tor}_n^R(M,N)
\longrightarrow\operatorname{Tor}_n^R(M,N'')\\
&\xrightarrow{\delta}\operatorname{Tor}_{n-1}^R(M,N')
\longrightarrow\cdots\qquad(n\ge1).
\end{aligned}
\]
\end{theorem}
\begin{proof}
在第一变量取马蹄消解，得到逐次分裂的短正合列 \(0\to P_n'\to P_n\to P_n''\to0\)。加性函子 \(-\otimes_RN\) 保持分裂短正合列，故张量后仍为复形短正合列，可以取同调长正合列。这里保证左端单射的是分裂性；张量仅有右正合性并不足以得到它。第二变量可用平衡性改在该侧取马蹄消解，或固定右模投射消解并用各 \(P_n\) 平坦。由\cref{b4:thm:tor-balanced}证明中的连接态射相容性，两种构造对应同一长正合列；零次识别给出所列张量尾端。
\end{proof}

\begin{example}\label{b4:exm:tor-cyclic}
设 \(R\) 是主理想整环，\(0\ne a\in R\)，\(N\) 是幺性 \(R\) 模。将 \(R/(a)\) 的两项消解与 \(N\) 张量，得到链次数一、零上的 \(N\xrightarrow{\times a}N\)，故零次同调是乘 \(a\) 的余核，一次同调是其核：
\[
\operatorname{Tor}_0^R(R/(a),N)=N/aN,
\qquad \operatorname{Tor}_1^R(R/(a),N)=\{\,n\in N\mid an=0\,\},
\]
且高次为零。对 \(R=\mathbb Z\)、\(N=\mathbb Z/n\mathbb Z\)、\(a=m\)，其中 \(m,n>0\)，一次 Tor 同构于 \(\mathbb Z/\gcd(m,n)\mathbb Z\)。与\cref{b4:exm:ext-cyclic}比较，对同一个乘法映射 \(a:N\to N\)，\(\operatorname{Ext}_R^1(R/(a),N)\) 取其余核 \(N/aN\)，\(\operatorname{Tor}_1^R(R/(a),N)\) 则取其核。有限循环群算例给出同构的交换群，并不使两者成为同一构造。
\end{example}

设 \(k\) 为域，令 \(R=k[x,y]\)，并把 \(k\) 看作商模 \(R/(x,y)\)。\cref{b4:prop:koszul-resolution}给出的两变量消解中，\(P_0=R\) 记录一个生成元，\(P_1=R^2\) 的微分为 \((a,b)\mapsto xa+yb\)，记录两个关系 \(x,y\)，而 \(P_2=R\) 的微分为 \(c\mapsto(-yc,xc)\)，记录这两个关系之间的关系 \(x(-y)+yx=0\)。与 \(k\) 张量后，\(x,y\) 都作用为零，所以微分全为零，各项 \(k,k^2,k\) 直接成为同调：
\[
\operatorname{Tor}_0^R(k,k)\cong k,\qquad
\operatorname{Tor}_1^R(k,k)\cong k^2,\qquad
\operatorname{Tor}_2^R(k,k)\cong k,
\]
更高次为零。这里一次和二次 Tor 的维数分别记录了这个消解中的两个关系以及它们之间的一条关系。

% !TeX root = ../../Book4.tex
% 纲要标题：### 9.4 一次 Ext、Tor 与正合性判据
% 纲要标签：b4:sec:0804
% 纲要位置：计划纲要.md，第 3420 行。

\section{一次 Ext、Tor 与正合性判据}
\label{b4:sec:0804}

一次 Ext 或 Tor 的消失能检测投射性、内射性或平坦性，但测试对象必须遍历相应的全部模。整数模与多项式理想的例子可以把正合性的失败具体算出来。


% 纲要内容（仅作写作计划，尚非教材正文）：
%
% * Ext 检验投射性和内射性
% * Tor 检验平坦性
% * 回到第二卷反例解释失败机制
%

\subsection{Ext 与投射性、内射性}
\label{b4:subsec:0804:01}

\begin{theorem}\label{b4:thm:ext-projective-injective}
设 \(R\) 是含幺环，\(P,I\) 是幺性左 \(R\) 模。下列判据成立：
\begin{enumerate}
\item \(P\) 投射，当且仅当对所有幺性左 \(R\) 模 \(N\)，\(\operatorname{Ext}_R^1(P,N)=0\)；这也等价于对所有这样的 \(N\) 及所有 \(n>0\)，\(\operatorname{Ext}_R^n(P,N)=0\)。
\item \(I\) 内射，当且仅当对所有幺性左 \(R\) 模 \(M\)，\(\operatorname{Ext}_R^1(M,I)=0\)；这也等价于对所有这样的 \(M\) 及所有 \(n>0\)，\(\operatorname{Ext}_R^n(M,I)=0\)。
\end{enumerate}
\end{theorem}
\begin{proof}
投射或内射时用集中于零次的相应消解，正次 Ext 消失。反过来，设对所有左模 \(N\) 都有 \(\operatorname{Ext}_R^1(P,N)=0\)。任取满射 \(q:B\to C\)，令 \(A=\ker q\)，第二变量长正合列给出
\[
\operatorname{Hom}_R(P,B)\xrightarrow{q\circ-}\operatorname{Hom}_R(P,C)
\longrightarrow\operatorname{Ext}_R^1(P,A)=0.
\]
所以每个 \(P\to C\) 都是某个 \(P\to B\) 与 \(q\) 的复合，即 \(P\) 有投射提升性质。

同样，设对所有左模 \(M\) 都有 \(\operatorname{Ext}_R^1(M,I)=0\)。任取单射 \(j:A\to B\)，令 \(C=\operatorname{coker}j\)，第一变量长正合列给出
\[
\operatorname{Hom}_R(B,I)\xrightarrow{-\circ j}\operatorname{Hom}_R(A,I)
\longrightarrow\operatorname{Ext}_R^1(C,I)=0.
\]
所以每个 \(A\to I\) 都能延拓为 \(B\to I\)，即 \(I\) 内射。
\end{proof}

“对所有模”的量词有实质作用。例如 \(\operatorname{Ext}_{\mathbb Z}^1(\mathbb Z/2\mathbb Z,\mathbb Q)=0\)，并不能推出 \(\mathbb Z/2\mathbb Z\) 投射：这次消失来自第二变量 \(\mathbb Q\) 的内射性，而换成第二变量 \(\mathbb Z\) 时，Ext 就是 \(\mathbb Z/2\mathbb Z\)。又如 \(\operatorname{Ext}_{\mathbb Z}^1(\mathbb Z,\mathbb Z)=0\)，不能推出第二个 \(\mathbb Z\) 内射：这次消失来自第一变量 \(\mathbb Z\) 的投射性，换成第一变量 \(\mathbb Z/2\mathbb Z\) 就得到同一个非零 Ext。检验投射性时固定第一变量、遍历第二变量；检验内射性时则相反。

\subsection{Tor 与平坦性}
\label{b4:subsec:0804:02}

\begin{theorem}\label{b4:thm:tor-flat-criterion}
设 \(R\) 是含幺环，\(N\) 是幺性左 \(R\) 模。下列条件等价：\(N\) 平坦；对每个幺性右 \(R\) 模 \(M\)，\(\operatorname{Tor}_1^R(M,N)=0\)；对每个这样的 \(M\) 及每个 \(n>0\)，\(\operatorname{Tor}_n^R(M,N)=0\)。右模的判据以相反侧的幺性左模作测试。
\end{theorem}
\begin{proof}
若 \(N\) 平坦，\(-\otimes_RN\) 正合，故右模投射消解在张量后仍在正次数正合；于是所有正次 Tor 为零。全消失当然蕴含一次消失。若一次均消失，对任意右模单射 \(A\to B\)，令 \(C=B/A\)。长正合列的片段
\[
\operatorname{Tor}_1^R(C,N)\longrightarrow A\otimes_RN
\longrightarrow B\otimes_RN
\]
表明后一个箭头单射。对任意右模短正合列 \(0\to A\to B\to C\to0\)，张量积的右正合性已经保证 \(A\otimes_RN\to B\otimes_RN\to C\otimes_RN\to0\) 正合；缺少的只是最左边的单射性。现在这一步也成立，故 \(-\otimes_RN\) 正合，即 \(N\) 平坦。反环给出右模版本。
\end{proof}

不必为了判定平坦性计算所有高次 Tor；一次已经足够。相反，只对某一个模计算一次 Tor，通常只能发现反例，不能得出整体平坦性。

\subsection{非正合性的模论反例}
\label{b4:subsec:0804:03}

设 \(m>1\) 为整数。把 \(0\to\mathbb Z\xrightarrow{\times m}\mathbb Z\to\mathbb Z/m\mathbb Z\to0\) 与 \(\mathbb Z/m\mathbb Z\) 作张量，乘 \(m\) 的映射变为零。由于 \(\mathbb Z\) 投射，相应长正合列的片段是
\[
0\longrightarrow\operatorname{Tor}_1^{\mathbb Z}(\mathbb Z/m\mathbb Z,\mathbb Z/m\mathbb Z)
\longrightarrow\mathbb Z/m\mathbb Z\xrightarrow{0}\mathbb Z/m\mathbb Z.
\]
零映射的核是整个 \(\mathbb Z/m\mathbb Z\)，故 \(\operatorname{Tor}_1^{\mathbb Z}(\mathbb Z/m\mathbb Z,\mathbb Z/m\mathbb Z)\cong\mathbb Z/m\mathbb Z\)，它记录了张量后丢失的单射性。对同一列施加 \(\operatorname{Hom}_{\mathbb Z}(-,\mathbb Z)\)，得到
\[
\mathbb Z\xrightarrow{\times m}\mathbb Z
\longrightarrow\operatorname{Ext}_{\mathbb Z}^1(\mathbb Z/m\mathbb Z,\mathbb Z)
\longrightarrow0.
\]
乘 \(m\) 的映射不是满射，其余核就是 \(\operatorname{Ext}_{\mathbb Z}^1(\mathbb Z/m\mathbb Z,\mathbb Z)\cong\mathbb Z/m\mathbb Z\)。对同一个整数短正合列，张量积丢失的单射性由 Tor 的核项记录，反变 Hom 丢失的满射性由 Ext 的余核项记录。

\ProseParagraphBreak
再取域 \(k\) 上的多项式环 \(R=k[x,y]\) 及其理想 \(\mathfrak m=(x,y)\)，将 \(k\) 识别为 \(R/\mathfrak m\)。因为 \(R\) 交换，以下将这些模同时视为右 \(R\) 模。短正合列 \(0\to\mathfrak m\to R\to k\to0\) 在第一变量的 Tor 长正合列中给出
\[
\underbrace{\operatorname{Tor}_2^R(R,k)}_{=0}\longrightarrow
\operatorname{Tor}_2^R(k,k)\longrightarrow\operatorname{Tor}_1^R(\mathfrak m,k)
\longrightarrow\underbrace{\operatorname{Tor}_1^R(R,k)}_{=0}.
\]
两端因 \(R\) 投射而消失，再代入上一节算出的二次 Tor，得到
\[
\operatorname{Tor}_1^R(\mathfrak m,k)\cong
\operatorname{Tor}_2^R(k,k)\cong k.
\]
故 \(\mathfrak m\) 不平坦，尽管它是整环中的无挠子模。这具体说明主理想整环上的无挠模平坦这一结论不能直接推广到任意整环。上面的连接态射已经将二次 Tor 与核模的一次 Tor 对应；下一节说明何时能这样降低次数。

% !TeX root = ../../Book4.tex
% 纲要标题：### 9.5 维数平移与合冲模初步
% 纲要标签：b4:sec:0805
% 纲要位置：计划纲要.md，第 3426 行。

\section{维数平移与合冲模初步}
\label{b4:sec:0805}

短正合列中的投射或内射项消去高次导出群，使相邻次数之间出现自然同构。逐次取合冲模或余合冲模，便把高次计算移到更低次数。


% 纲要内容（仅作写作计划，尚非教材正文）：
%
% * 对短正合列 \(0\to\Omega M\to P\to M\to0\) 及投射模 \(P\)，由长正合列推出正次数 Ext 和 Tor 的维数平移，并注明变量侧别
% * 对内射嵌入 \(0\to N\to I\to C\to0\) 给出对偶的 Ext 平移公式；区分合冲模选择与导出不变量的自然同构
% * 迭代平移与截短消解，说明在有限步后出现投射合冲模的含义；不在此预设一般同调维数理论
% * 本节是第三卷第15–21章所用同调工具包的一部分；第11.1–11.2节在其基础上统一定义和研究各类同调维数
%
% ---
%

\subsection{投射消解中的 Ext 与 Tor 维数平移}
\label{b4:subsec:0805:01}

\begin{definition}\label{b4:def:syzygy-basic}
设 \(R\) 是含幺环，\(M\) 是幺性左 \(R\) 模。给定投射满射 \(P\to M\)，其核记为 \(\Omega M\)，称为与这一选择对应的第一\term[hechongmo]{合冲模}[syzygy module]。在固定投射消解 \(P_\bullet\xrightarrow{\varepsilon}M\) 中，令 \(\Omega^0M=M\)、\(\pi_0=\varepsilon:P_0\to M\)，并递归定义
\[
\Omega^{r+1}M=\ker\pi_r,\qquad \pi_r:P_r\twoheadrightarrow\Omega^rM.
\]
这里 \(r\ge1\) 时，消解的正合性给出 \(\operatorname{im}\mathrm{d}_r=\Omega^rM\subseteq P_{r-1}\)，所以微分 \(\mathrm{d}_r:P_r\to P_{r-1}\) 唯一因子化为满射 \(\pi_r\) 与包含映射的复合。这就确定了递归中使用的映射。右模用右投射消解作同一定义。
\end{definition}
\symidx{\(\Omega^rM\)：相对于选定投射消解的第 \(r\) 次合冲模}

\begin{theorem}[基础维数平移]\label{b4:thm:dimension-shift-basic}
设 \(R\) 是含幺环，\(0\to K\to P\to M\to0\) 是幺性左 \(R\) 模短正合列，且 \(P\) 投射。则对任意幺性左 \(R\) 模 \(N\)、幺性右 \(R\) 模 \(A\) 及整数 \(i\ge1\)，连接态射给出同构
\[
\operatorname{Ext}_R^i(K,N)\cong\operatorname{Ext}_R^{i+1}(M,N),
\qquad
\operatorname{Tor}_{i+1}^R(A,M)\cong\operatorname{Tor}_i^R(A,K).
\]
Ext 同构对 \(N\) 及所给短正合列之间的态射自然；Tor 同构对 \(A\) 及所给短正合列之间的态射自然。低次则为正合列
\[
\begin{aligned}
0&\longrightarrow\operatorname{Hom}_R(M,N)\longrightarrow\operatorname{Hom}_R(P,N)
\longrightarrow\operatorname{Hom}_R(K,N)\\
&\longrightarrow\operatorname{Ext}_R^1(M,N)\longrightarrow0,\\
0&\longrightarrow\operatorname{Tor}_1^R(A,M)\longrightarrow A\otimes_RK
\longrightarrow A\otimes_RP\\
&\longrightarrow A\otimes_RM\longrightarrow0.
\end{aligned}
\]
若原列为幺性右 \(R\) 模短正合列，则 Tor 结论改为 \(\operatorname{Tor}_{i+1}^R(M,N)\cong\operatorname{Tor}_i^R(K,N)\)，其中 \(N\) 为幺性左 \(R\) 模；该同构对 \(N\) 及所给右模短正合列之间的态射自然。
\end{theorem}
\begin{proof}
第一变量 Ext 长正合列在相关位置为
\[
\operatorname{Ext}_R^i(P,N)\longrightarrow\operatorname{Ext}_R^i(K,N)
\longrightarrow\operatorname{Ext}_R^{i+1}(M,N)
\longrightarrow\operatorname{Ext}_R^{i+1}(P,N).
\]
\(i\ge1\) 时两端由投射性消失。Tor 的相应片段是
\[
\operatorname{Tor}_{i+1}^R(A,P)\longrightarrow\operatorname{Tor}_{i+1}^R(A,M)
\longrightarrow\operatorname{Tor}_i^R(A,K)\longrightarrow\operatorname{Tor}_i^R(A,P),
\]
两端由投射蕴含平坦而消失。当 \(i=0\) 时，Ext 片段左端是 \(\operatorname{Hom}_R(P,N)\)，Tor 片段右端是 \(A\otimes_RP\)，它们一般不为零，因此连接态射未必是同构；保留这些零次项才得到所列低次正合列。自然性来自原长正合列对短正合列态射的自然性；右模情形把同一证明用于另一变量。
\end{proof}

这一操作称为\term[weishu-pingyi]{维数平移}[dimension shifting]。一条短正合列只使次数降低一步；逐次取合冲模并重复这一操作，才能把高次问题降到一次。零次仍须使用包含 Hom 或张量项的正合列。

\subsection{内射消解中的 Ext 平移与自然性}
\label{b4:subsec:0805:02}

设 \(R\) 是含幺环，\(N\) 是幺性左 \(R\) 模。给定到内射模的单射 \(j:N\to I\)，其余核 \(C=I/j(N)\) 称为与这一选择对应的第一\term[yuhechongmo]{余合冲模}[cosyzygy module]。在内射消解 \(0\to N\to I^0\xrightarrow{\mathrm{d}^0}I^1\xrightarrow{\mathrm{d}^1}\cdots\) 中，第一次取商给出
\[
C^1=I^0/N\cong\operatorname{im}\mathrm{d}^0=\ker \mathrm{d}^1\subseteq I^1.
\]
因此 \(\mathrm{d}^0\) 在这个商上诱导单射 \(C^1\to I^1\)，再取 \(C^2=I^1/C^1\)，就有短正合列 \(0\to C^1\to I^1\to C^2\to0\)。这里将 \(N\) 与其在 \(I^0\) 中的像识别，并将 \(C^1\) 与上述单射的像识别。不断取这样的余核，与投射消解中不断取核相对偶。

\begin{proposition}\label{b4:prop:injective-dimension-shift}
设 \(R\) 是含幺环，\(0\to N\to I\to C\to0\) 为幺性左 \(R\) 模短正合列，\(I\) 内射。对每个幺性左 \(R\) 模 \(M\) 和整数 \(i\ge1\)，有同构
\[
\operatorname{Ext}_R^i(M,C)\cong\operatorname{Ext}_R^{i+1}(M,N).
\]
该同构对 \(M\) 及所给短正合列之间的态射自然。零次给出
\[
\operatorname{Ext}_R^1(M,N)\cong
\operatorname{coker}\bigl(\operatorname{Hom}_R(M,I)\to\operatorname{Hom}_R(M,C)\bigr).
\]
\end{proposition}
\begin{proof}
第二变量长正合列在正次的片段为
\[
\operatorname{Ext}_R^i(M,I)\longrightarrow\operatorname{Ext}_R^i(M,C)
\longrightarrow\operatorname{Ext}_R^{i+1}(M,N)
\longrightarrow\operatorname{Ext}_R^{i+1}(M,I).
\]
\(i\ge1\) 时两端由内射性消失，故中间连接态射为同构；零次则留下
\[
\begin{aligned}
0&\longrightarrow\operatorname{Hom}_R(M,N)\longrightarrow\operatorname{Hom}_R(M,I)
\longrightarrow\operatorname{Hom}_R(M,C)\\
&\longrightarrow\operatorname{Ext}_R^1(M,N)\longrightarrow0,
\end{aligned}
\]
从而得到所述余核表达。连接态射的自然性来自长正合列的自然性。
\end{proof}

余合冲模 \(C\) 取决于所选内射嵌入，正如 \(\Omega M\) 取决于所选投射满射。这里的自然性所比较的是给定短正合列及其态射；更换消解时，比较定理提供相应的比较映射。

\begin{lemma}[Schanuel 引理]\label{b4:lem:schanuel}
设 \(R\) 是含幺环，\(0\to K\to P\to M\to0\)、\(0\to K'\to P'\to M\to0\) 为两个幺性左 \(R\) 模短正合列。若中间模 \(P,P'\) 都投射，则 \(K\oplus P'\cong K'\oplus P\)。
\end{lemma}
\begin{proof}
记两个满射为 \(\pi:P\to M\)、\(\pi':P'\to M\)，取纤维积
\[
E=P\times_M P'=\{(p,p')\in P\oplus P':\pi(p)=\pi'(p')\}.
\]
任取 \(p'\in P'\)，由 \(\pi\) 满射，可取 \(p\in P\) 使 \(\pi(p)=\pi'(p')\)，于是 \((p,p')\in E\) 映到 \(p'\)。故第二投影 \(E\to P'\) 满射，其核为 \(\{(p,0):\pi(p)=0\}\cong K\)；同理第一投影 \(E\to P\) 满射，其核为 \(\{(0,p'):\pi'(p')=0\}\cong K'\)。这样得到短正合列
\[
0\longrightarrow K\longrightarrow E\longrightarrow P'\longrightarrow0,
\qquad
0\longrightarrow K'\longrightarrow E\longrightarrow P\longrightarrow0.
\]
因右端投射，两列均分裂，故 \(E\cong K\oplus P'\cong K'\oplus P\)。
\end{proof}

这个\term[schanuel-yinli]{Schanuel 引理}[Schanuel's lemma]说明两种合冲模在加上投射直和因子后同构。\footnote{沙努埃尔（Stephen Hoel Schanuel，1933-2014），美国数学家，研究同调代数、范畴论与数论。}这不能推出 \(K\cong K'\)，两者甚至可能不同构。例如在非零环 \(R\) 上，自由模 \(M\) 可以取恒等满射，核为零；也可以取投影 \(M\oplus R\to M\)，核为非零的 \(R\)。此外，两条短正合列的分裂一般不唯一，所得直和同构也依赖于分裂的选择。

\subsection{迭代平移与截短消解}
\label{b4:subsec:0805:03}

\begin{proposition}\label{b4:prop:iterated-dimension-shift}
设 \(R\) 是含幺环，\(M,N\) 是幺性左 \(R\) 模，\(A\) 是幺性右 \(R\) 模。固定 \(M\) 的投射消解 \(P_\bullet\to M\) 及其 \(r\) 次合冲模 \(\Omega^rM\)，其中 \(r\ge0\) 为整数。则对整数 \(i\ge1\)，有
\[
\operatorname{Ext}_R^i(\Omega^rM,N)\cong\operatorname{Ext}_R^{i+r}(M,N),
\qquad
\operatorname{Tor}_i^R(A,\Omega^rM)\cong\operatorname{Tor}_{i+r}^R(A,M).
\]
若固定 \(N\) 的内射消解 \(N\to I^\bullet\)，令 \(C^0=N\)、\(C^{j+1}=\operatorname{coker}(C^j\to I^j)\)（\(j\ge0\)）。这里 \(C^0\to I^0\) 为增广单射；对 \(j\ge1\)，内射消解的正合性使 \(I^{j-1}\to I^j\) 在前一步的商 \(C^j\) 上诱导单射，其像为 \(\ker(I^j\to I^{j+1})\)。则
\[
\operatorname{Ext}_R^i(M,C^r)\cong\operatorname{Ext}_R^{i+r}(M,N).
\]
投射平移的两个同构对另一变量及所给投射消解之间的增广链映射自然；内射平移的同构对 \(M\) 及所给内射消解之间的增广上链映射自然。此外，固定消解中的 \(\Omega^rM\) 投射，当且仅当对所有幺性左 \(R\) 模 \(N\)，\(\operatorname{Ext}_R^{r+1}(M,N)=0\)。
\end{proposition}
\begin{proof}
对各短正合列 \(0\to\Omega^{j+1}M\to P_j\to\Omega^jM\to0\) 依次用\cref{b4:thm:dimension-shift-basic}。例如两步给出
\[
\operatorname{Ext}_R^{i+2}(M,N)\cong\operatorname{Ext}_R^{i+1}(\Omega M,N)
\cong\operatorname{Ext}_R^i(\Omega^2M,N).
\]
每次将次数降低一，同时将第一变量换成下一次合冲模；复合 \(r\) 次连接同构即得第一式。Tor 用相同的短正合列迭代，内射式则对 \(0\to C^j\to I^j\to C^{j+1}\to0\) 反复用\cref{b4:prop:injective-dimension-shift}。增广链映射在各核上诱导映射，增广上链映射在各余核上诱导映射；这些映射组成短正合列之间的态射，故每一步的连接同构及其复合都自然。

最后，对每个幺性左 \(R\) 模 \(N\)，第一式在 \(i=1\) 时成为
\[
\operatorname{Ext}_R^1(\Omega^rM,N)\cong\operatorname{Ext}_R^{r+1}(M,N).
\]
因此右边对所有这样的 \(N\) 消失，当且仅当左边也对所有 \(N\) 消失，而\cref{b4:thm:ext-projective-injective}说明后者当且仅当 \(\Omega^rM\) 投射。
\end{proof}

若 \(r=0\)，这一条件就是 \(M\) 本身投射。若 \(r\ge1\) 且 \(\Omega^rM\) 投射，原消解可截成
\[
0\longrightarrow\Omega^rM\longrightarrow P_{r-1}\longrightarrow\cdots
\longrightarrow P_0\longrightarrow M\longrightarrow0.
\]
这里把 \(\Omega^rM\) 作为第 \(r\) 次投射项，所谓长度不超过 \(r\)，指投射项的最高非零次数不超过 \(r\)，不把增广目标 \(M\) 计作投射项。例如长度零表示 \(M\) 投射；长度不超过一表示存在 \(0\to P_1\to P_0\to M\to0\)，其中 \(P_0,P_1\) 均投射。反过来，若有长度不超过 \(r\) 的投射消解，施加 \(\operatorname{Hom}_R(-,N)\) 后，第 \(r+1\) 次项为零，所以对每个 \(N\) 都有 \(\operatorname{Ext}_R^{r+1}(M,N)=0\)。由上述判据，任意投射消解的第 \(r\) 次合冲模都投射。高次 Ext 的全称消失与投射消解能在第 \(r\) 次截短因此等价；\chapref{b4:ch:10}再用它定义并比较各种同调维数。

\subsection{维数平移在局部代数中的应用}
\label{b4:subsec:0805:04}

第三卷的局部代数以各项有限秩的自由消解记录生成元和关系。这里已经得到其中所需的基本同调推理：把 \(0\to\Omega M\to P\to M\to0\) 送入 Ext 或 Tor 的长正合列，利用 \(P\) 的正次导出函子消失，把次数降低一步。有限性和局部性在选择极小自由消解、利用 Nakayama 引理\footnote{中山正（假名：なかやま ただし，罗马音：Tadashi Nakayama，1912-1964），日本数学家，研究环论与表示论。}时另有作用；上述维数平移本身对任意含幺环成立。

\begin{example}\label{b4:exm:local-shift-koszul}
设 \(k\) 是域，\(R=k[x,y]_{(x,y)}\)，\(\mathfrak m=(x,y)R\)，剩余域记为 \(k\)。由于 \(R\) 交换，以下各左模也通过 \(m\cdot r=r\cdot m\) 视为右 \(R\) 模。\(R\) 是整环，\(R/(x)\cong k[y]_{(y)}\) 也是整环，所以 \(x,y\) 满足\cref{b4:prop:koszul-resolution}的假设，给出自由消解
\[
0\longrightarrow R\xrightarrow{\binom{-y}{x}}R^2
\xrightarrow{(x\;y)}R\longrightarrow k\longrightarrow0.
\]
与 \(k\) 张量后，两个微分映射均为零，故 \(\operatorname{Tor}_2^R(k,k)\cong k\)。短正合列 \(0\to\mathfrak m\to R\to k\to0\) 遂给出 \(\operatorname{Tor}_1^R(\mathfrak m,k)\cong k\ne0\)。由 Tor 的平坦判据，\(\mathfrak m\) 不平坦；投射模必平坦，故它也不投射。

\ProseParagraphBreak
另一方面，在这条消解中，\(\Omega k=\operatorname{im}(R^2\to R)=\mathfrak m\)，而正合性给出
\[
\Omega^2k=\ker\bigl(R^2\xrightarrow{(x\;y)}R\bigr)
=\operatorname{im}\bigl(R\xrightarrow{\binom{-y}{x}}R^2\bigr)\cong R.
\]
最后的同构使用了最左边微分的单射性，所以 \(\Omega^2k\) 投射，而 \(\Omega k\) 不投射。这同时给出长度二的自由消解和阻止在第一步截短的一次 Tor。
\end{example}


\begin{chaptersummary}
导出函子先由消解定义，比较定理保证选择独立，马蹄引理给出长正合列，消去性进一步确定它的泛性质。Ext 对第一变量反变、第二变量协变；Tor 对两个变量协变且要求相反的模侧别。平衡性由同一个双复形的两种计算证明，因此可以选择较容易处理的变量作消解。

\ProseParagraphBreak
一次 Ext 与 Tor 足以分别检验投射、内射和平坦性，但测试变量必须遍历所有相应模。维数平移利用投射或内射项的高次消失，把高次问题逐步移到合冲模或商对象。零次仍含 Hom 或张量项，不能套用正次数同构。这些工具既用于下一章的扩张解释，也用于后续同调维数及第三卷的局部自由消解。
\end{chaptersummary}

\BookExercises

\begin{exercise}\label{b4:ex:08-derived-exactness}
设 \(\mathcal A\) 有足够多内射对象，\(F:\mathcal A\to\mathcal B\) 为交换范畴间加性左正合函子。证明 \(F\) 正合当且仅当 \(R^1F(A)=0\) 对每个 \(A\) 成立，并说明此时所有正次均消失。
\end{exercise}
\begin{solution}
若 \(F\) 正合，内射消解的增广正合性经 \(F\) 保留，故所有正次上同调为零。若一次全消失，对任意 \(0\to A\to B\to C\to0\)，长正合列使 \(F(B)\to F(C)\) 满；结合已有左正合性，\(F\) 正合。最后再用正向结论即得全消失。
\end{solution}

\begin{exercise}\label{b4:ex:08-derived-natural-transformation}
设 \(F,G:\mathcal A\to\mathcal B\) 是加性左正合函子，\(\mathcal A\) 有足够多内射对象，\(a:F\to G\) 自然。用同一内射消解构造 \(R^na:R^nF\to R^nG\)，并证明它与泛性给出的延拓相同。
\end{exercise}
\begin{solution}
对 \(A\to I_A^\bullet\)，各分量 \(a_{I_A^q}\) 因自然性与微分交换，故形成复形映射。取上同调，利用比较映射的自然性得到自然变换 \(R^na\)。对马蹄消解施加 \(a\) 给出短正合复形列之间的映射，故它与连接态射相容。零次通过核的识别正是 \(a_A\)，所以由\cref{b4:thm:derived-universal}的唯一性，等于泛延拓。
\end{solution}

\begin{exercise}\label{b4:ex:08-central-actions}
设 \(R\) 是交换环，\(M,N\) 为 \(R\) 模。证明 \(r\in R\) 通过 \(M\) 的乘法映射或 \(N\) 的乘法映射诱导的 Ext 自同态相同。对 Tor 也证明相应结论。
\end{exercise}
\begin{solution}
对 \(M\) 的投射消解，逐次乘 \(r\) 是乘法态射的提升。在 \(\operatorname{Hom}_R(P_n,N)\) 上，预合成 \(r\) 与后合成 \(r\) 相同，因为 \(f(rp)=r\cdot f(p)\)。因此在上同调上相同。Tor 的消解张量中有 \(rp\otimes n=p\otimes rn\)，故同样成立。若 \(rM=0\) 或 \(rN=0\)，对应零态射诱导零，因此该 \(r\) 零化所有相关 Ext 与 Tor。
\end{solution}

\begin{exercise}\label{b4:ex:08-ext-sum-product}
设 \(R\) 是含幺环，\((M_i)_{i\in I}\)、\(N\) 是左模。证明对每个 \(n\ge0\)，有
\[
\operatorname{Ext}_R^n(\bigoplus_iM_i,N)\cong\prod_i\operatorname{Ext}_R^n(M_i,N).
\]
指出无限指标处所用的集合论选择。
\end{exercise}
\begin{solution}
为各 \(M_i\) 选择投射消解后，逐次直和仍正合。其每项也投射：给定满射 \(X\twoheadrightarrow Y\) 及 \(\bigoplus_iP_i\to Y\)，先为每个 \(P_i\to Y\) 选择一个到 \(X\) 的提升，再由直和的泛性质合成提升。无限指标时，这里要从一族非空的提升集合中作选择。因此这些消解的直和是 \(\bigoplus_iM_i\) 的投射消解，Hom 将它变为各 Hom 复形的直积。

交换群的直积中核逐分量计算；直积也保持满射，但这是另一处选择：对一族满射 \(X_i\twoheadrightarrow Y_i\) 和给定的 \((y_i)\)，为每个 \(y_i\) 选择原像 \(x_i\)，才能组成积中的原像。因而积保持短正合列，复形的圈与边界也均可逐分量计算，得到
\[
H^n\!\left(\prod_i C_i^\bullet\right)
\cong\prod_i H^n(C_i^\bullet).
\]
将此式用于上述 Hom 复形即得结论。
\end{solution}

\begin{exercise}\label{b4:ex:08-ext-map-cyclic}
设 \(m,n\) 为正整数且 \(n\mid m\)，取商映射 \(q:\mathbb Z/m\mathbb Z\to\mathbb Z/n\mathbb Z\)。对交换群 \(N\)，在 \(\operatorname{Ext}_{\mathbb Z}^1(\mathbb Z/r\mathbb Z,N)\cong N/rN\) 的识别下，求 \(q^*\)。
\end{exercise}
\begin{solution}
分别用 \(\mathbb Z\xrightarrow{m}\mathbb Z\) 和 \(\mathbb Z\xrightarrow{n}\mathbb Z\) 作源、目标的自由消解。商映射 \(q\) 在零次提升为 \(f_0=\operatorname{id}_{\mathbb Z}\)，一次提升为 \(f_1=(m/n)\operatorname{id}_{\mathbb Z}\)；链映射条件正是
\[
n\circ f_1=f_0\circ m,\qquad n(m/n)=m.
\]
对这条链映射作 Hom，箭头反向，一次映射由预合成 \(f_1\) 给出，即乘 \(m/n\)。因此
\[
q^*:N/nN\longrightarrow N/mN,\qquad [a]\longmapsto[(m/n)a].
\]
若 \(a\) 改变 \(nb\)，像改变 \(mb\)，故良定义。这个公式表明 Ext 的反变方向不是简单地沿自然取商。
\end{solution}

\begin{exercise}\label{b4:ex:08-tor-map-cyclic}
仍设 \(m,n\) 为正整数且 \(n\mid m\)，对交换群 \(N\) 求上述 \(q\) 诱导的 \(\operatorname{Tor}_1^{\mathbb Z}(\mathbb Z/m\mathbb Z,N)\to\operatorname{Tor}_1^{\mathbb Z}(\mathbb Z/n\mathbb Z,N)\)。
\end{exercise}
\begin{solution}
沿前一题的链映射作张量，箭头保持原方向，一次分量仍为乘 \(m/n\)。在一次同调即乘法映射的核上，诱导映射为
\[
\{\,a\in N\mid ma=0\,\}\longrightarrow\{\,b\in N\mid nb=0\,\},
\qquad a\longmapsto(m/n)a.
\]
\(n(m/n)a=ma=0\)，确实进入目标。同一条消解链映射在这里给出 \(N[m]\to N[n]\)，而前一题中给出 \(N/nN\to N/mN\)：方向的区别来自张量的协变性与 Hom 第一变量的反变性。
\end{solution}

\begin{exercise}\label{b4:ex:08-integer-four-six}
计算 \(\operatorname{Hom}_{\mathbb Z}(\mathbb Z/4\mathbb Z,\mathbb Z/6\mathbb Z)\)、\(\operatorname{Ext}_{\mathbb Z}^1(\mathbb Z/4\mathbb Z,\mathbb Z/6\mathbb Z)\)、\(\operatorname{Tor}_1^{\mathbb Z}(\mathbb Z/4\mathbb Z,\mathbb Z/6\mathbb Z)\)，并分别写出具体代表。
\end{exercise}
\begin{solution}
三者的抽象群均为 \(\mathbb Z/2\mathbb Z\)。Hom 由 \(\bar1\) 的像 \(0\) 或 \(3\) 决定；Tor 是 \(\mathbb Z/6\mathbb Z\) 中乘四的核 \(\{0,3\}\)；Ext 是 \((\mathbb Z/6\mathbb Z)/4(\mathbb Z/6\mathbb Z)\)，其子群 \(4(\mathbb Z/6\mathbb Z)=\{0,2,4\}\)，可用 \(0,1\) 代表两个类。核中的元素和余核中的类承担不同角色。
\end{solution}

\begin{exercise}\label{b4:ex:08-tor-conormal}
设 \(R\) 是交换环，\(I\) 是理想。证明有典范同构 \(\operatorname{Tor}_1^R(R/I,R/I)\cong I/I^2\)，无需假定 \(I\) 平坦。
\end{exercise}
\begin{solution}
对 \(0\to I\to R\to R/I\to0\) 在另一变量取 \(R/I\)，长正合列及 \(R\) 投射给出
\[
0\longrightarrow\operatorname{Tor}_1^R(R/I,R/I)
\longrightarrow I\otimes_RR/I\longrightarrow R\otimes_RR/I.
\]
最后一个映射在纯张量上将 \(a\otimes\bar r\)（\(a\in I\)）送到 \(R\otimes_RR/I\) 中的同名张量，而
\[
a\otimes\bar r=1\otimes\overline{ar}=0,
\]
故整个映射为零。张量商同构 \(I\otimes_RR/I\to I/I^2\) 将 \(a\otimes\bar r\) 送到 \(ar+I^2\)，所以连接态射与此同构的复合给出所求典范同构。
\end{solution}

\begin{exercise}\label{b4:ex:08-truncated-cubic-ring}
设 \(k\) 是域，令 \(R=k[t]/(t^3)\)，并将 \(k\) 识别为 \(R/(t)\)。写出 \(k\) 的周期自由消解，计算 \(\operatorname{Ext}_R^n(k,k)\)、\(\operatorname{Tor}_n^R(k,k)\) 以及 \(\operatorname{Ext}_R^n(k,R)\) 的正次数。
\end{exercise}
\begin{solution}
取各 \(P_i=R\)，消解为
\[
\cdots\xrightarrow{t}R\xrightarrow{t^2}R
\xrightarrow{t}R\xrightarrow{t^2}R\xrightarrow{t}R
\longrightarrow k\longrightarrow0.
\]
从 \(P_1\to P_0\) 起，微分交替为乘 \(t\)、乘 \(t^2\)。由于
\[
\begin{aligned}
\ker(t:R\to R)&=(t^2)=\operatorname{im}(t^2:R\to R),\\
\ker(t^2:R\to R)&=(t)=\operatorname{im}(t:R\to R),
\end{aligned}
\]
每个正次数的位置均正合，零次的核则为 \((t)\)。以 \(k\) 作 Hom 目标或张量因子时，\(t\) 的作用为零，故微分全为零，得到
\[
\operatorname{Ext}_R^n(k,k)\cong k,\qquad
\operatorname{Tor}_n^R(k,k)\cong k\quad(n\ge0).
\]
以 \(R\) 作 Hom 目标时，上链微分仍交替为乘 \(t\)、乘 \(t^2\)，上述核像等式使所有正次 Ext 为零；零次是 \(\ker(t:R\to R)=(t^2)\cong k\)。
\end{solution}

\begin{exercise}\label{b4:ex:08-one-module-test-fails}
给出一个非投射模 \(M\) 与一个非零模 \(N\)，使 \(\operatorname{Ext}^n(M,N)=0\) 对全部 \(n>0\) 成立；再给出非内射 \(N\) 和非零 \(M\) 的相应例子。
\end{exercise}
\begin{solution}
在 \(\mathbb Z\) 上取 \(M=\mathbb Z/2\mathbb Z\)、\(N=\mathbb Q\)，第二变量内射使正次全消失，而 \(M\) 不是投射。再取 \(M=N=\mathbb Z\)，第一变量投射使正次全消失，\(N\) 却不内射，因为 \(2\mathbb Z\to\mathbb Z\)、\(2a\mapsto a\) 无法延拓。两个例子都说明判据中的全称测试不能缩成任意一次测试。
\end{solution}

\begin{exercise}\label{b4:ex:08-flat-kernel}
设 \(R\) 是含幺环，\(0\to A\to B\to C\to0\) 为左模短正合列，\(B,C\) 平坦。证明 \(A\) 平坦。
\end{exercise}
\begin{solution}
对任意右模 \(M\)，长正合列含有 \(\operatorname{Tor}_2^R(M,C)\to\operatorname{Tor}_1^R(M,A)\to\operatorname{Tor}_1^R(M,B)\)。两端因平坦而为零，故中间为零。再由\cref{b4:thm:tor-flat-criterion}得到 \(A\) 平坦。
\end{solution}

\begin{exercise}\label{b4:ex:08-extension-flat}
设 \(0\to A\to B\to C\to0\) 是任意含幺环上的左模短正合列，\(A,C\) 平坦。证明 \(B\) 平坦。若 \(A,B\) 平坦，结论能否改为 \(C\) 平坦？
\end{exercise}
\begin{solution}
对任意右模 \(M\)，一次 Tor 片段 \(\operatorname{Tor}_1(M,A)\to\operatorname{Tor}_1(M,B)\to\operatorname{Tor}_1(M,C)\) 两端为零，故 \(B\) 平坦。反向断言不成立：\(0\to\mathbb Z\xrightarrow2\mathbb Z\to\mathbb Z/2\mathbb Z\to0\) 前两项自由，末项与自身的一次 Tor 非零，故不平坦。
\end{solution}

\begin{exercise}\label{b4:ex:08-zero-degree-shift}
以 \(0\to2\mathbb Z\to\mathbb Z\to\mathbb Z/2\mathbb Z\to0\) 为例，说明不能把 \(\operatorname{Ext}^i(\Omega M,N)\cong\operatorname{Ext}^{i+1}(M,N)\) 延伸到 \(i=0\)。写出正确的余核公式。
\end{exercise}
\begin{solution}
取 \(N=\mathbb Z\)。分别以 \(f\mapsto f(1)\) 和 \(g\mapsto g(2)\) 识别 \(\operatorname{Hom}(\mathbb Z,\mathbb Z)\) 与 \(\operatorname{Hom}(2\mathbb Z,\mathbb Z)\)，两者均为 \(\mathbb Z\)。限制映射将 \(f(1)=a\) 的同态送到 \(g(2)=f(2)=2a\) 的同态，因此在这些坐标下为乘二。长正合列给出的正确公式是
\[
\operatorname{Ext}^1_{\mathbb Z}(\mathbb Z/2\mathbb Z,\mathbb Z)
\cong\operatorname{coker}\!\left(
\operatorname{Hom}(\mathbb Z,\mathbb Z)\longrightarrow
\operatorname{Hom}(2\mathbb Z,\mathbb Z)\right)
\cong\mathbb Z/2\mathbb Z.
\]
它不是整个 \(\operatorname{Hom}(2\mathbb Z,\mathbb Z)\)；必须先模去能够从 \(\mathbb Z\) 延拓来的同态。
\end{solution}

\begin{exercise}\label{b4:ex:08-injective-shift-integers}
设 \(m\) 为正整数。用 \(0\to\mathbb Z\to\mathbb Q\to\mathbb Q/\mathbb Z\to0\) 和内射方向的低次平移，重新求 \(\operatorname{Ext}_{\mathbb Z}^1(\mathbb Z/m\mathbb Z,\mathbb Z)\)，并给出同构的具体元素。
\end{exercise}
\begin{solution}
\(\operatorname{Hom}(\mathbb Z/m\mathbb Z,\mathbb Q)=0\)，故低次正合列给出 \(\operatorname{Ext}^1(\mathbb Z/m\mathbb Z,\mathbb Z)\cong\operatorname{Hom}(\mathbb Z/m\mathbb Z,\mathbb Q/\mathbb Z)\)。一个同态由 \(\bar1\) 的像决定，该像必须被 \(m\) 消去。若 \(z+\mathbb Z\in\mathbb Q/\mathbb Z\) 满足 \(mz\in\mathbb Z\)，写 \(a=mz\)，就有 \(z=a/m\)；反之，每个 \(a/m+\mathbb Z\) 都被 \(m\) 消去。两个整数 \(a,b\) 给出同一像，当且仅当 \(m\mid(a-b)\)。因此
\[
\mathbb Z/m\mathbb Z\longrightarrow\operatorname{Hom}(\mathbb Z/m\mathbb Z,\mathbb Q/\mathbb Z),
\qquad
\bar a\longmapsto\bigl(\bar1\mapsto a/m+\mathbb Z\bigr)
\]
是同构，再由连接态射将这些同态对应到 Ext 类。
\end{solution}

\begin{exercise}\label{b4:ex:08-syzygy-change}
设 \(R\) 为含幺环，\(P,M\) 为幺性左 \(R\)-模，\(P\) 投射。给定满射 \(p:P\twoheadrightarrow M\)，记 \(K=\ker p\)，并取任意投射左模 \(Q\)。构造另一个从投射模到 \(M\) 的满射，使其核同构于 \(K\oplus Q\)。证明对任意左 \(R\)-模 \(N\)、右 \(R\)-模 \(A\) 及 \(n\ge1\)，有
\[
\operatorname{Ext}_R^n(K\oplus Q,N)\cong\operatorname{Ext}_R^n(K,N),\qquad
\operatorname{Tor}_n^R(A,K\oplus Q)\cong\operatorname{Tor}_n^R(A,K).
\]
\end{exercise}
\begin{solution}
取 \(P\oplus Q\to M\)、\((x,q)\mapsto p(x)\)，其核为 \(K\oplus Q\)。有限双积的投射消解可逐项作双积，Hom 或张量后仍分成两部分。\(Q\) 的正次导出函子消失，故所得正次 Ext、Tor 分别等于 \(K\) 的相应值。零次却保留额外项：
\[
\begin{aligned}
\operatorname{Hom}_R(K\oplus Q,N)
&\cong\operatorname{Hom}_R(K,N)\oplus\operatorname{Hom}_R(Q,N),\\
A\otimes_R(K\oplus Q)
&\cong(A\otimes_RK)\oplus(A\otimes_RQ).
\end{aligned}
\]
这些额外项一般非零，因而正次数的同构不能直接延伸到零次。
\end{solution}

\begin{exercise}\label{b4:ex:08-truncate-independent}
设 \(R\) 为含幺环，\(M\) 为幺性左 \(R\)-模，且 \(M\) 有长度不超过 \(r\) 的投射消解，其中 \(r\ge1\)。证明在任意另一投射消解中，第 \(r\) 次合冲模投射，而第 \(r-1\) 次未必投射。
\end{exercise}
\begin{solution}
已知短消解使 \(\operatorname{Ext}_R^{r+1}(M,N)=0\) 对全部 \(N\) 成立。另一消解的迭代平移给出 \(\operatorname{Ext}_R^1(\Omega^rM,N)=0\)，故该合冲模投射。取域 \(k\)，令 \(R=k[x,y]\)、\(M=R/(x,y)\)、\(r=2\)。Koszul 自由消解的二次项为 \(R\)，而其一次合冲模为 \((x,y)\)。张量 \(k\) 后两个微分均为零，故 \(\operatorname{Tor}_2^R(k,k)\cong k\)；正次数平移给出 \(\operatorname{Tor}_1^R((x,y),k)\cong k\ne0\)。因此 \((x,y)\) 不平坦，从而不投射；二次合冲模却同构于 \(R\)，是投射模。
\end{solution}
